\BeleskeID{09_1-s024}
% element: 09_1-s024:n01
Početna ravnoteža struja u drugoj oblasti glasi
\[\frac{V_{DD}-V_O}{R_L}=\frac{k_n}{2}\frac{L_nE_{Cn}}{L_nE_{Cn}+(V_I-V_{Tn})}(V_I-V_{Tn})^2\]
U tački odlučivanja ulazni i izlazni napon su jednaki. Posle zamene $V_I=V_O=V_S$, ulazna nepoznata pojavljuje se i u brojicu i u imeniocu. Da bi se izdvojio kvadratni izraz po prekoračenju praga, pišemo
\[\frac{V_{DD}-(V_S-V_{Tn})-V_{Tn}}{R_L}=\frac{k_n}{2}\frac{L_nE_{Cn}}{L_nE_{Cn}+(V_S-V_{Tn})}(V_S-V_{Tn})^2\]
Množenjem imeniteljem i sređivanjem dobija se prva kvadratna jednačina na slajdu. Deljenjem koeficijentom uz $(V_S-V_{Tn})^2$ dobija se druga, normalizovana jednačina. Zatim se unose brojne vrednosti i bira radna tačka u pretpostavljenoj oblasti.

\[\begin{aligned}\left(\frac{k_nL_nE_{Cn}}2+\frac1{R_L}\right)(V_S-V_{Tn})^2+\left(\frac{L_nE_{Cn}}{R_L}-\frac{V_{DD}-V_{Tn}}{R_L}\right)(V_S-V_{Tn})-L_nE_{Cn}\frac{V_{DD}-V_{Tn}}{R_L}&=\\&\quad 0\end{aligned}\]

